\documentclass[conference]{IEEEtran}
\usepackage{graphicx,amsmath,hyperref,natbib,booktabs}

\title{GNOmE: Zero-Query Mechanistic Interpretability via Computation Graph Reading}

\author{
\IEEEauthorblockN{Sehaj Singh}
\IEEEauthorblockA{Independent Research\\
sehajr-singhs@gmail.com}
}

\begin{document}
\maketitle

\begin{abstract}
We introduce GNOmE, a method for mechanistic interpretability that extracts a neural network's computation graph from Jacobian-weighted contributions in a single forward pass, then reads the graph with a GNN to predict per-unit importance without model interventions. On 6 trained transformers across IOI and Induction tasks, GNOmE correlates with ground-truth importance at $r = 0.748$ (vs.\ path patching $r = -0.365$), achieves 0.96 cross-task transfer, and reduces query complexity from $O(n)$ to $O(1)$.
\end{abstract}

\begin{IEEEkeywords}
mechanistic interpretability, graph neural networks, circuit analysis, transformer interpretability
\end{IEEEkeywords}

\section{Introduction}

Mechanistic interpretability decomposes neural network computation into subgraph circuits \citep{olah2020zoom, elhage2021mathematical}. Path patching \citep{wang2023interpretability} and activation patching \citep{geiger2021causal} probe these circuits by zeroing individual units, requiring $O(n_\text{layers} \times n_\text{heads})$ forward passes. We show this can be reduced to $O(1)$ by reading the model's intrinsic computation graph.

\section{Method}

\textbf{Graph Extraction.} For each attention head and MLP layer, we compute its contribution vector as the mean output across batch and positions. Edges between consecutive layers are cosine similarities between these vectors. Node features encode layer depth, unit type (head vs.\ MLP), and graph connectivity (in/out degree).

\textbf{GNN Reader.} A 2-layer message-passing GNN takes the graph as input and predicts per-node importance. Once trained on a few models, it generalizes to unseen models without accessing them.

\section{Experiments}

We train 2-layer 4-head transformers (105K params) on IOI and Induction tasks (3 seeds each). Table~\ref{tab:results} summarizes results.

\begin{table}[t]
\centering
\caption{Comparison with ground-truth zero-ablation importance.}
\label{tab:results}
\begin{tabular}{lccc}
\toprule
Method & Corr.\ ($r$) & P@3 & Query Cost \\
\midrule
GNOmE & \textbf{0.748} & \textbf{0.667} & $O(1)$ \\
Path Patching & $-0.365$ & 0.278 & $O(n_L n_H)$ \\
\bottomrule
\end{tabular}
\end{table}

\textbf{Cross-Model Transfer.} GNN trained on IOI predicts Induction importance at $r = 0.954$, and vice versa at $r = 0.963$. LOO cross-validation: $r = 0.864 \pm 0.201$.

\textbf{Threshold Sweep.} Optimal at $\tau = 0.3$: $r = 0.965$ with 24\% fewer edges, demonstrating robust pruning.

\section{Conclusion}

GNOmE reduces mechanistic interpretability from $O(n)$ interventions to $O(1)$ graph extractions, outperforming path patching while enabling scalability to large models.

\bibliographystyle{IEEEtran}
\bibliography{references}

\end{document}
